\section{Tails, deletion, and limits of positive comparison methods}
\label{p1:limits:section}

The stopped expansion controls a distribution without imposing an
exponential moment bound on the original count.  This distinction is
necessary in the present model.  We first record the exact support and
local structure, then give a nearby count with factorial tails.  A real
arithmetic cylinder shows why these tails cannot simply be transferred
back.  The final parts examine deletion of active pairs and the positive
cost of forcing private prime signs.

\subsection{Exact support and the local rank-one structure}

Write $Z=Z_N$, $p=2^{-L}$ and $R_N=\lceil N/L\rceil$.
For $j\ge1$, $(z)_j=z(z-1)\cdots(z-j+1)$.

\begin{proposition}[Support and factorial lattice]
\label{p1:limits:lattice}
For every $N\ge2$ and $L\ge1$,
\begin{equation}
 0\le Z\le R_N,\qquad
 \E(Z)_j\in j!p^j\mathbb Z_{\ge0}\quad(j\ge1).
 \label{p1:limits:lattice-eq}
\end{equation}
\end{proposition}
\begin{proof}
Two starts at distance less than $L$ cannot both be active: the
predecessor of the later start lies inside the earlier constant run.
Thus the active starts in an interval of $N$ integers are $L$-separated,
which gives the support bound.  For $j$ distinct starts, their $jL$
relative sign prescriptions form an affine binary system.  Its joint
probability is either zero or $2^{-r}$ for an integer $r\le jL$.
It is therefore an integer multiple of $p^j$.  Finally,
$(Z)_j=j!\sum_{|I|=j}\prod_{x\in I}J_x$.
\end{proof}

The local exceptional windows have an exact description. We use the
rank-one conclusion of Lemma~\ref{p1:limits:local-rigidity}; a bound $o(L)$
for the local nullity would not suffice for what follows.



The proof of Lemma~\ref{p1:limits:local-rigidity} also gives raw-value nullity at most two at critical length,
because the even relation code has codimension at most one in the
full relation code.  It also gives exact $k$-wise independence of the star coordinates
$Y_i=f(U)f(U+i)$, $1\le i<B$, whenever
$(k+1)\log(64B)<\log U$: every nonempty product of at most $k$
relative signs selects at most $k+1$ values, and Runge excludes a
square product.  These are sufficient bounds.  They do not establish
that critical windows with a nonzero relation actually exist.

Let $\mathcal R$ be the collection of nonempty even sets $S$ of integers
whose product is a square and which are contained in at least one
$V_x$, $x\in I_N$.  Write
\[
 \begin{gathered}
 a_S=\min S,\qquad b_S=\max S,\qquad d_S=b_S-a_S,\\
 B_S=\{b_S-L+1,\ldots,a_S+1\}\cap I_N,\qquad
 e_S=\ind_{\{a_S+1\in I_N\}}.
 \end{gathered}
\]

\begin{proposition}[Exact local blocks]
\label{p1:limits:local-blocks}
For sufficiently large $N$ in a fixed critical band, the intervals
$B_S$, $S\in\mathcal R$, are nonempty and pairwise disjoint.
They partition the locally deficient starts.  On such a block,
\begin{equation}
 \P(J_x=1)=
 \begin{cases}0,&x=a_S+1,\\2p,&x\in B_S\setminus\{a_S+1\},\end{cases}
 \qquad
 \sum_{x\in B_S}J_x\sim\mathrm{Bernoulli}\bigl(2p(|B_S|-e_S)\bigr).
 \label{p1:limits:block-law}
\end{equation}
Outside these blocks the marginal is $p$.  In particular,
\begin{align}
 \E Z-Np&=p\sum_{S\in\mathcal R}(|B_S|-2e_S),
       \label{p1:limits:block-mean}\\
 \sum_{x\in I_N}|\P(J_x=1)-p|
       &=p\sum_{S\in\mathcal R}|B_S|.
       \label{p1:limits:block-absolute}
\end{align}
A full block contributes $p(L-d_S-1)$ to the mean discrepancy.
There are at most two blocks truncated by the endpoints of $I_N$,
whose total boundary correction is at most $2Lp$ in absolute value.
\end{proposition}
\begin{proof}
The condition $S\subseteq V_x$ is exactly
$b_S-L+1\le x\le a_S+1$.  If two distinct such sets belonged
to the same window, they would give two independent nonzero
relative relations, contrary to Lemma~\ref{p1:limits:local-rigidity}.
Thus their intervals are disjoint and every deficient window lies
in exactly one.  Its unique nonzero relation has affine phase
$(-1)^{\ind_{\{x-1\in S\}}}$, so Fourier inversion gives
$p(1+(-1)^{\ind_{\{x-1\in S\}}})$.  Since $x-1\le a_S$,
the predecessor belongs to $S$ exactly when $x=a_S+1$.
This proves the marginal formulas.  The diameter of $B_S$ is at
most $L-1$; consequently at most one of its starts is active.
Its count is Bernoulli with mean the sum of the displayed
marginals.  Summation proves \eqref{p1:limits:block-mean} and
\eqref{p1:limits:block-absolute}.  A full interval has
$L-d_S+1$ starts.  Disjoint intervals allow at most one truncation
at each endpoint, and each full interval has at most $L$ starts.
\end{proof}

No independence between the Bernoulli block counts is asserted.
For the complete relative sign vector $Y_x\in\{\pm1\}^L$ one
also has the exact identity
\[
 \TV(\Law Y_x,\mathrm{Unif}\{\pm1\}^L)=1-2^{-r_x},
 \qquad r_x=L-\dim S_x.
\]
Indeed $Y_x$ is uniform on a binary image of size $2^{L-r_x}$.
Thus this local distance is either zero or $1/2$ at critical
length.  A nonzero $S$ must have diameter at least
$(1-o(1))\log N$, by the subcritical part of Lemma~\ref{p1:limits:local-rigidity};
its full block has length at most
$(1/\log2-1+o(1))\log N$.

\subsection{A nearby count with factorial tails at every order}

Let $G=G_N$, $n=|G|$, $\mu=np$, and let $E'_N$ be the exceptional
event of Lemma~\ref{p1:core:exceptional}.  It contains the
occurrence of a primitive family of size at least three, an active
contact, an active bad dependent pair, two good active pairs, and
an additional dependence after compressing the unique good pair.
Its probability $\delta_N$ is $O_\varepsilon(N^{-4/5+\varepsilon})$.
Outside it, the active original spaces are a direct sum of single
windows and at most one good pair.  Deleting one endpoint of that
pair, if present, therefore gives a direct family.

\begin{theorem}[A modified count with all factorial tails]
\label{p1:limits:modified}
On the original probability space define
$\widetilde Z_N=Z_N^\circ\ind_{(E'_N)^c}$.
Then
\begin{equation}
 \TV(\Law Z_N,\Law\widetilde Z_N)
       \ll_\varepsilon N^{-4/5+\varepsilon}.
 \label{p1:limits:modified-tv}
\end{equation}
Uniformly for every integer $k\ge1$, every real $s\ge1$, and
every integer $r\ge1$,
\begin{align}
 \P(Z_N^\circ\ge k)&\le\frac{\mu^{k-1}}{(k-1)!}+\delta_N,
       \label{p1:limits:real-tail}\\
 \P(\widetilde Z_N\ge k)&\le\frac{\mu^{k-1}}{(k-1)!},
       \label{p1:limits:modified-tail}\\
 \E s^{\widetilde Z_N}&\le s(1+p(s-1))^n
                         \le s e^{(s-1)\mu},
       \label{p1:limits:modified-exponential}\\
 \E(\widetilde Z_N)_r&\le\mu^r+r\mu^{r-1}.
       \label{p1:limits:modified-moments}
\end{align}
\end{theorem}
\begin{proof}
Let $R$ be the maximum cardinality of a direct active subfamily of
$G$.  Directness makes the prescribed events exactly independent.
A union bound gives
\[
 \P(R\ge k)\le\sum_{\substack{|I|=k\\I\ \text{direct}}}q(I)
       \le\binom nkp^k\le\frac{\mu^k}{k!}.
\]
Choose one maximizing direct family in each realization.  Its
$\binom Rr$ subsets of size $r$ are all direct and active, whence
\[
 \E\binom Rr\le\binom nrp^r\quad(r\ge0).
\]
Outside $E'_N$ we have $Z_N^\circ\le R+1$, and everywhere
$\widetilde Z_N\le R+1$.  The two tail bounds follow.
Expanding $s^R=(1+(s-1))^R$ with nonnegative coefficients gives
$\E s^R\le(1+p(s-1))^n$.  Finally
$\binom{R+1}{r}=\binom Rr+\binom R{r-1}$ proves the factorial
moment bound.  The coupling error is at most $\delta_N$ on $G$;
Proposition~\ref{p1:arith:core} restores the deleted starts.
\end{proof}

These inequalities concern the modified variable.  The event defining
it depends on the whole active configuration, so its indicators no
longer have the original affine marginals or centred factorization.
In particular, \eqref{p1:limits:modified-tv} does not bound
$\E[s^{Z_N};E'_N]$.

\subsection{Effective tails of the original count}

A different argument gives a quantitative tail bound for $Z_N$ itself,
at a larger threshold.  It tracks the moment order explicitly.
Put $H=\log N$ and $h=\log H$.

\begin{lemma}[Branching with an explicit order dependence]
\label{p1:limits:branching}
For any $j$ old starts, let $\mathcal F_{\rm old}$ contain all
their raw sign values.  There is a deterministic exceptional set of
new starts of size at most
\[
 D_{j,N}=j(2L+1)+B\sqrt{3N}
       \exp\left(C\sqrt{m_{j,N}/\log(m_{j,N}+2)}\right),
 \qquad
 m_{j,N}=\left\lceil\pi(B)+\frac{jB\log(3N)}{\log B}\right\rceil,
\]
outside which the $L$ new relative signs are uniform and independent
of $\mathcal F_{\rm old}$.  In particular, uniformly for $1\le j\le H$,
\begin{equation}
 \log(\lambda_N+D_{j,N})
       \le\frac H2+C_C\left(h+\frac{\sqrt j\,H}{h}\right).
 \label{p1:limits:branching-uniform}
\end{equation}
If $F_k=\E(Z_N)_k$, then
\begin{equation}
 F_k\le C_C\exp\left\{
   \frac{k-2}{2}H+C_C\left(kh+\frac{k^{3/2}H}{h}\right)
                     \right\}\quad(2\le k\le H).
 \label{p1:limits:growing-moments}
\end{equation}
\end{lemma}
\begin{proof}
Let $\mathcal P$ consist of the primes at most $B$ and the primes
larger than $B$ dividing one of the old values.  It has at most
$m_{j,N}$ members.  The number of integers at most $3N$ whose
squarefree part is supported on $\mathcal P$ is at most
\[
 \sqrt{3N}\prod_{\ell\in\mathcal P}(1+\ell^{-1/2}).
\]
For a set of $m$ primes, the estimate
$\ell_i\gg i\log(i+1)$ for its $i$th member implies
\[
 \log\prod_{\ell\in\mathcal P}(1+\ell^{-1/2})
       \le C\sqrt{m/\log(m+2)}.
\]
Delete starts whose new window overlaps an old one or contains
at least two integers with squarefree part supported on
$\mathcal P$.  Counting each such integer in at most $B$ windows
gives the displayed bound for $D_{j,N}$.

If a nontrivial combination of the new relative rows belonged
to the old raw-value span, it would select an even nonempty set
of new values.  A prime larger than $B$ in such a value is private
in the new window, so must also divide an old value.  Every
selected value would consequently have squarefree part supported
on $\mathcal P$.  At least two are selected, contradicting the
deletion rule.  Thus the rank increases by exactly $L$ even after
conditioning on all old raw values, which proves independence.

For $j\le H$, $m_{j,N}=O_C(jH^2/h)$ and
$\log(m_{j,N}+2)\gg h$.  This proves
\eqref{p1:limits:branching-uniform}.  Summing the conditional
probabilities of the next start, using the bound one on the
exceptional set, yields the finite recurrence
$F_{j+1}\le(\lambda_N+D_{j,N})F_j$.
The pair estimate gives $F_2=O_C(1)$.  Multiplying from $j=2$
to $k-1$ and summing $\sqrt j$ proves
\eqref{p1:limits:growing-moments}.
\end{proof}

\begin{proposition}[A quantified effective support]
\label{p1:limits:effective-tail}
For each fixed $D>0$ and integer $r\ge0$, there is a constant
$K=K(C,D,r)$ such that
\begin{equation}
 \E[Z_N^r;Z_N>A_N]\le N^{-D},\qquad
 A_N=\left\lceil\sqrt N\exp\left(\frac{K\log N}
                                    {(\log\log N)^{2/3}}\right)\right\rceil
 \label{p1:limits:effective-threshold}
\end{equation}
for sufficiently large $N$.  Hence, for every fixed
$\eta,D>0$ and integer $r\ge0$,
$\E[Z_N^r;Z_N>N^{1/2+\eta}]=O(N^{-D})$.
\end{proposition}
\begin{proof}
Take $k=\lceil h^{2/3}\rceil$.  Eventually $k>r$ and
$A_N\ge2k$.  On $Z_N>A_N$,
$Z_N^r\le2^kA_N^{r-k}(Z_N)_k$.
The logarithm of the resulting bound from
\eqref{p1:limits:growing-moments} is at most
\[
 \frac{r-2}{2}H+C_C\left(kh+\frac{k^{3/2}H}{h}\right)
        -(k-r)\frac{KH}{h^{2/3}}+k\log2+O_C(1).
\]
Here $kh=o(H)$, $k^{3/2}/h=1+o(1)$ and
$(k-r)/h^{2/3}=1+o(1)$.  A sufficiently large fixed $K$
makes this at most $-DH$.  The last assertion follows by
comparison of thresholds.
\end{proof}

The all-fixed-order bounds of Theorem~\ref{p1:fixed:moments}
also imply $\E[Z_N^r;Z_N>N^\eta]=O(N^{-D})$ for fixed
$\eta,D,r$, by choosing one sufficiently large fixed moment.
The value of \eqref{p1:limits:effective-threshold} is its explicit
control of a growing order, without extrapolating fixed-order
constants.

A larger threshold admits a much smaller exceptional probability.
The distinction is useful when the moment order itself increases
as fast as $\log N$.

\begin{proposition}[A very small tail beyond a near-linear threshold]
\label{p1:limits:colored-tail}
Let $S=\sqrt{\log N\log\log N}$ and fix
$0<c<\sqrt{13}/4$.  For
$A_{N,c}=\lceil R_Ne^{-cS}\rceil$, one has, for large $N$,
\begin{equation}
 \P(Z_N>A_{N,c})\le e^{-(\log N)^3},\qquad
 \E[Z_N^r;Z_N>A_{N,c}]\le e^{-(\log N)^3/2}
       \quad(1\le r\le\log N).
 \label{p1:limits:colored-tail-bound}
\end{equation}
Using only the published estimate of de la Bret\`eche--Drappeau,
the same conclusion holds for $c<2/\sqrt5$.
\end{proposition}
\begin{proof}
Write $\kappa(n)=P^+(\sf(n))$, with $P^+(1)=1$.
We first give the analytic input and its transfer to these odd
prime pivots.  Pascadi's bound \cite[Corollary~7.2]{pascadi2025},
with the forms $n,n+h$, implies, for fixed $\eta>0$ and
$\gamma=13/8-\eta$,
\begin{equation}
 \sum_{1\le h\le B}
 \#\{n\in[N-B,2N+B]:\kappa(n),\kappa(n+h)\le e^{aS}\}
 \le N\left\{
 e^{-(\gamma/(2a)+o(1))S}
       +e^{-(a+1/(2a)+o(1))S}\right\},
 \label{p1:limits:two-small-pivots}
\end{equation}
uniformly for $a$ in a fixed compact subinterval of $(0,\infty)$.
Here the parameter in the cited result is fixed before $N$:
its smoothness range contains $y=e^{aS}$ and its coefficient
range contains every $h\le B=O(\log N)$.
For ordinary $y$-smooth pairs the estimate follows from
$\Psi(X,y)=X\varrho(\log X/\log y)e^{O(\log X/\log y)}$
and $\log\varrho(u)=-(1+o(1))u\log u$.

The difference between ordinary smoothness and
$\kappa(n)\le y$ must also be counted.  Such a difference
requires $n=\ell^2m$ for a prime $\ell>y$ and
$\kappa(m)\le y$.  For $s>1/2$, positive summation and
the exact Dirichlet series give
\[
 \#\{n\le4N:\kappa(n)\le y<P^+(n)\}
 \le(4N)^s\zeta(2s)\prod_{\ell\le y}(1+\ell^{-s})
                          \sum_{\ell>y}\ell^{-2s}.
\]
Put $w=\log y=aS$, $u=\log N/w$ and $s=1-(\log u)/w$.
Chebyshev's prime bound and partial summation show that the
logarithm of the Euler product is
$O(\log w+u/\log u)=o(S)$.  Extending the last sum to all
integers gives $O(y^{1-2s}/(2s-1))$.  The resulting upper bound
is $N e^{-(a+1/(2a)+o(1))S}$.  Each exceptional integer lies
in at most $2B$ relevant pairs, proving
\eqref{p1:limits:two-small-pivots}.  Theorem~4.1 of
\cite{bretedrappeau2020} gives the same argument with
$\gamma=8/5-\eta$.

Choose $\eta>0$ so small that
$a=\sqrt{\gamma/2}>c$, and set $Y=e^{aS}$.
Delete windows with at least two pivots at most $Y$.
By \eqref{p1:limits:two-small-pivots}, their number is
$N e^{-(a+o(1))S}$; the second term there is strictly smaller.
Conditional on all prime signs at most $Y$, each remaining
window has marginal $p$: retain its possible defective value
as anchor and use a private prime greater than $Y$ for every
other value.  Join two remaining windows when they share a
prime greater than $Y$ with odd valuation.  This is an exact
conditional dependency graph.  A greedy coloring uses at most
\[
 \chi=\left\lceil1+\frac{B^2\log(3N)}{\log Y}(N/Y+2)\right\rceil
\]
colors.  Conditional independence in each color class and
H\"older's inequality imply
\[
 \E[e^{tZ_G}\mid\mathcal F_Y]
       \le\exp\{\lambda_N(e^{\chi t}-1)/\chi\},\qquad t\ge0.
\]
Chernoff's bound, with $\chi t=\log(z/\lambda_N)$, is at most
$\exp\{-(z/\chi)(\log(z/\lambda_N)-1)\}$.
Take $z=2\chi(\log N)^2$.  Since
$\log\chi=\log N-aS+O(\log\log N)$, this probability is
at most $e^{-(\log N)^3}$.  The deleted count plus $z$ is
$o(R_Ne^{-cS})$, proving the first assertion.  Finally
$Z_N\le R_N\le N$ proves the moment assertion uniformly for
$r\le\log N$.
\end{proof}

\subsection{An arithmetic cylinder with large exponential moments}

\begin{proposition}[Divergence of the real exponential moment]
\label{p1:limits:cylinder}
Along odd primes $q\to\infty$, put $N_q=2^{q-1}$ and
$L_q=q-1$.  For every fixed $t>2$,
\begin{equation}
 \liminf_{q\to\infty}\frac{\log N_q}{N_q}
        \log\E e^{tZ_{N_q,L_q}}\ge(t-2)\log2>0.
 \label{p1:limits:exponential-divergence}
\end{equation}
The same bound holds for the count of positive runs alone.
\end{proposition}
\begin{proof}
Let $M=2N_q+q$ and prescribe the prime-sign cylinder
\[
 E_q=\{f(q)=-1,\ f(\ell)=1
             \text{ for every prime }\ell\le M,\ \ell\ne q\}.
\]
It has probability $2^{-\pi(M)}$.  On this event,
$f(n)=(-1)^{v_q(n)}$ for all integers used by the windows.
Every start $x=uq+1\in I_{N_q}$ with $q\nmid u(u+1)$ has a
negative predecessor, then exactly $q-1$ positive values, and
a negative next value.  Their number is
\[
 K_q=(1-2/q)N_q/q+O(1)=(1-o(1))N_q/q.
\]
Therefore
$\log\E e^{tZ}\ge tK_q-\pi(2N_q+q)\log2$.
Since $q=\log N_q/\log2+1$, the prime number theorem gives
\eqref{p1:limits:exponential-divergence}.
\end{proof}

The number $2$ is a sufficient threshold in this construction;
its optimality is not asserted.  For every fixed $r$, the same
cylinder contributes at most $R_{N_q}^r2^{-\pi(M)}=o(1)$ to
the $r$th moment.  It is thus compatible with all fixed-moment
estimates and with the total-variation approximation.

The cylinder also produces connected families that lie far beyond
the bounded-nullity boundary used in the proof.  The next elementary
rank argument extracts one such family with a positive centred
coefficient.

\begin{lemma}[Deleting large private ranks]
\label{p1:limits:peeling}
Let $A$ be a compatible family of $m$ affine events, each with
constraint-space dimension $L$, and put
$R=\dim\sum_{x\in A}S_x$ and $p=2^{-L}$.
For an integer $1\le t\le L$, repeatedly remove a member whose
current private rank is at least $t$.  The final family $A'$ satisfies
\[
 |A'|\ge m-R/t,\qquad
 \dim\sum_{x\in A'}S_x\le R,\qquad
 r_x(A')<t\quad(x\in A'),
\]
where $r_x(I)=\dim\sum_{y\in I}S_y-
\dim\sum_{y\in I\setminus\{x\}}S_y$.
Each direct component $C$ of $A'$ satisfies
\begin{equation}
 c(C)\ge(1-|C|p2^t)q(C)\ge(1-mp2^t)q(C).
 \label{p1:limits:positive-core}
\end{equation}
If $m-R/t>0$ and $D=R/(m-R/t)$, some such component obeys
\begin{equation}
 \frac{\dim\sum_{x\in C}S_x}{|C|}\le D,\qquad
 |C|\ge L/D,\qquad \rho(C)\ge L^2/D-L.
 \label{p1:limits:rank-density}
\end{equation}
\end{lemma}
\begin{proof}
Each deletion reduces the current rank by at least $t$, so at
most $R/t$ deletions occur.  For a compatible family, all Fourier
phases of its relations are positive.  The centred Fourier formula is
\[
 c(C)=p^{|C|}\#\{(s_x)\in\ker(\bigoplus_{x\in C}S_x
                        \longrightarrow\sum_{x\in C}S_x):
                                      s_x\ne0\ (x\in C)\}.
\]
The image of the projection of this kernel on $S_x$ has dimension
$L-r_x(C)\ge L-t+1$.  In the uniform kernel, the probability
that $s_x=0$ is consequently at most $p2^t$.
A union bound and $q(C)=p^{|C|}2^{\rho(C)}$ prove
\eqref{p1:limits:positive-core}.  The size-weighted average of
the rank-to-size ratios of the components is at most $D$.
Choose a component no larger than that average.  Its rank is
at least $L$, proving $|C|\ge L/D$.  If $D<L$, use
$\rho(C)\ge|C|(L-D)$; if $D\ge L$, the stated lower bound
for nullity is nonpositive.
\end{proof}

\begin{proposition}[A real component of quadratic nullity]
\label{p1:limits:large-component}
Along the subsequence of Proposition~\ref{p1:limits:cylinder},
there are compatible connected families $C_q\subset G_{N_q}$
containing no good pair and satisfying
\begin{equation}
 \begin{gathered}
 |C_q|\ge(\kappa-o(1))L_q,\qquad
 \rho_*(C_q)=\rho(C_q)\ge(\kappa-o(1))L_q^2,\\
 c(C_q)\ge(1-O(L_q^{-1/2}))q(C_q)>0,\qquad
 \kappa=\frac{6\log2}{17}.
 \end{gathered}
 \label{p1:limits:quadratic-nullity}
\end{equation}
\end{proposition}
\begin{proof}
Restrict the starts in the cylinder to
\[
 T_q=\{x=uq+1:3N_q/2+4q\le x<2N_q,\quad q\nmid u(u+1)\},
 \qquad A_q=T_q\cap G_{N_q}.
\]
The regular-set estimate gives
$m=|A_q|\sim N_q/(2q)$; the cylinder witnesses compatibility.
For a good edge $x<y$ in channel $a:b$, its exact unit gives
$y=(b/a)x+(b/a)i-j$, with $-1\le i,j\le L_q-1$ and
$b/a\in\{4/3,3/2\}$.  The lower bound
$y\ge(b/a)x-(1+b/a)q$ exceeds $2N_q$ in this interval.
Thus $A_q$ has no good edge.

All prime coordinates used by these windows belong to the union
\[
 [1,2N_q/3+O(q)],\quad
 [3N_q/4-O(q),N_q+O(q)],\quad
 [3N_q/2-O(q),2N_q+O(q)].
\]
Indeed the multiplier of a prime in a used value is respectively
at least three, two, or one.  The prime number theorem bounds
their number by $(17/12+o(1))N_q/\log N_q$.
Hence the total rank $R$ satisfies
$R/m\le17/(6\log2)+o(1)$.
Apply Lemma~\ref{p1:limits:peeling} with
$t=\lfloor\frac12\log_2L_q\rfloor$.  Then $R/t=o(m)$,
$D\le17/(6\log2)+o(1)$ and $mp2^t=O(L_q^{-1/2})$.
Equations \eqref{p1:limits:positive-core}--\eqref{p1:limits:rank-density}
give the required component.  Every subfamily is compatible,
and nullity is nondecreasing under inclusion, so its compatible
nullity equals its full nullity.
\end{proof}

This is a statement about which arithmetic terms exist.  The
probability of the witnessing cylinder is not identified with
$q(C_q)$, and the proposition gives no lower bound for the total
variation error or for a sum of absolute component activities.

\subsection{Keeping positions while justifying pair deletion}
\label{p1:limits:positions}
\label{p14:comp:positional-proof}

Deleting an active pair is a transformation of the whole field.
The scalar approximation alone does not justify this operation.
We retain coordinates in the same finite expansion, using the
boundary and deficit estimates already proved.
Lemma~\ref{p1:limits:vector-reference} compares the unsigned
field \((J_x)_{x\in I_N}\) with independent Poisson single
arrivals and arrivals along good arithmetic pairs, with an
\(O_{C,j,\varepsilon}(N^{-4/5+\varepsilon})\) error for
every fixed polynomial weight of degree \(j\) in the total count.
The field retains integer positions, without run signs or exact
excess lengths; the target is on the same grid, not a diffuse
spatial process. We first work on the regular set \(G=G_N\)
and then restore the excluded starts.

On $\mathbb N^G$, let $\boldsymbol e_x$ denote the coordinate
unit vector and set
\[
 \Delta_x=\delta_{\boldsymbol e_x}-\delta_{\boldsymbol0},
 \qquad b_x=\delta_{\boldsymbol0}+p\Delta_x,
 \qquad \Delta_I=\mathop{*}_{x\in I}\Delta_x.
\]
For a hereditary class $\mathscr D$, the same partition by first
excluded prefix gives the exact identity
\begin{equation}
 \Law((J_x)_{x\in G})-
       \sum_{I\in\mathscr D}q(I)\Delta_I
 =\sum_{I\in\partial\mathscr D}\Delta_I*\boldsymbol\nu_I,
 \quad
 \boldsymbol\nu_I=
 \E\left[J_I\delta_{\sum_{y>\max I}J_y\boldsymbol e_y}\right].
 \label{p1:limits:vector-stop}
\end{equation}
One can obtain it directly by expanding
$\prod_x(1+u_x)$ with $u_x=J_x\Delta_x$ in the convolution
algebra.  The proof of Lemma~\ref{p1:stop:polynomial} permits
a different algebra element at every site.  Each
$\boldsymbol\nu_I$ is positive of mass $q(I)$, and
$\|\Delta_I\|_1=2^{|I|}$.  Thus the boundary estimate is
unchanged.

Use the field laws $\mathsf H_N$ and $\mathsf U_N$ and the intensities
$q_e,\iota_x,\nu_x$ defined in Section~\ref{p14:positions:section}. In
particular, $\beta=\sum_e q_e$, $\theta=2\beta$, $\lambda=Np$, and
$\iota_x=0$ outside $G$. The total-count laws are given by
\eqref{p1:limits:compound-target}.

\begin{proof}[Proof of Lemma~\ref{p1:limits:vector-reference}]
First work on $G$ in the unweighted norm.  The retained reference
is now
\[
 \mathsf B_\beta=\mathop{*}_{x\in G}b_x+
       \sum_{e=\{x,y\}\in E_N}q_e\Delta_x*\Delta_y*
                                      \mathop{*}_{z\in G\setminus e}b_z.
\]
The exact retained identity subtracts the missing direct
families, with coefficient $p^{|I|}$, and the missing compressed
extensions, with coefficient $q_ep^{|U|}$.  Taking norms of
these subtractions costs at most $F+\mathfrak G$, with exactly
the deficits of \eqref{p1:close:total-deficits}.  There is no
scalar cancellation in that estimate.  The vector version of
the first-boundary remainder is bounded by the same positive
sum as \eqref{p1:stop:norm}.  Hence
\[
 \|\Law((J_x)_{x\in G})-\mathsf B_\beta\|_1
       \ll_\varepsilon N^{-4/5+\varepsilon}.
\]

Put $\mathsf P=\exp_*\{p\sum_{x\in G}\Delta_x\}$.
Replacing the product of Bernoulli laws by $\mathsf P$ costs
$O(np^2)$.  In each edge term, replacing the product outside
$e$ by $\mathsf P$ costs $O(np^2+p)$ before multiplication
by $4q_e$.  Thus
\[
 \mathsf B_\beta=
 \mathsf P*\left(\delta_{\boldsymbol0}
                     +\sum_{e=\{x,y\}}q_e\Delta_x*\Delta_y\right)
              +O_{\ell^1}(np^2+\beta np^2+\beta p).
\]
The exponential remainder has norm at most
$e^{4\beta}-1-4\beta=O(\beta^2)$.  Moreover
\[
 p\sum_x\Delta_x+\sum_{e=\{x,y\}}q_e\Delta_x*\Delta_y
 =\sum_x(p-\iota_x)\Delta_x+
       \sum_{e=\{x,y\}}q_e
                    (\delta_{\boldsymbol e_x+\boldsymbol e_y}
                                      -\delta_{\boldsymbol0}),
\]
so this convolution exponential is exactly the positive law
$\mathsf H_N$ on $G$.  The errors $np^2$ and $\beta^2$ are
$O_\varepsilon(N^{-1+\varepsilon})$.  Restoring the deleted
starts in both laws costs $N^{-1+o(1)}$.

It remains to justify the weights; total variation alone is
insufficient.  Both total counts have every fixed moment
uniformly bounded: for the real count this is
Theorem~\ref{p1:fixed:moments}, and for the compound target
it follows from \eqref{p1:limits:compound-target} with bounded
parameters.  For their signed difference $\sigma$ and any fixed
integer $s>j$, H\"older's inequality gives
\[
 \sum(1+|\boldsymbol z|)^j|\sigma\{\boldsymbol z\}|
 \le\|\sigma\|_1^{1-j/s}
       \left(\sum(1+|\boldsymbol z|)^s
                         |\sigma\{\boldsymbol z\}|\right)^{j/s}.
\]
Choose $s$ sufficiently large after $\varepsilon$, and use the
unweighted estimate with a smaller error margin.  This proves
\eqref{p1:limits:vector-weighted} for each fixed $j$.
\end{proof}

The weights in this lemma have fixed polynomial degree.  They
are not exponential weights and do not contradict
Proposition~\ref{p1:limits:cylinder}.

Let $\Xi=\sum_{e\in E_N}J_e$, $C_N=\{\Xi\ge1\}$ and
$r_N=\P(C_N)$.  On $C_N$, choose one active good edge uniformly
and delete its two endpoints, obtaining a binary field $J^\flat$.
Off $C_N$ keep the field unchanged.  Also let $\widehat J$ be
the field obtained by deleting every vertex incident to an active
good edge.  Their counts are
\[
 W=\sum_xJ_x^\flat=Z-2\ind_{C_N},\qquad
 \widehat W=\sum_x\widehat J_x.
\]
The auxiliary random choice affects the field $J^\flat$ but
not $W$.

\begin{proposition}[Deletion, isolation, and remaining covariances]
\label{p1:limits:erasure}
Uniformly in the critical band, for every $\varepsilon>0$,
\begin{align}
 \TV(\Law J^\flat,\mathsf U_N)
       +\TV(\Law\widehat J,\mathsf U_N)
       &\ll_\varepsilon N^{-4/5+\varepsilon},
       \label{p1:limits:erasure-tv}\\
 \sum_x\big|\E J_x^\flat-(p-\iota_x)\big|
       +\sum_x\big|\E\widehat J_x-(p-\iota_x)\big|
       &\ll_\varepsilon N^{-1+\varepsilon}.
       \label{p1:limits:erasure-intensity}
\end{align}
For a counting field $X$ put
$K_{xy}(X)=\Cov(X_x,X_y)-\ind_{\{x=y\}}\E X_x$.
Then
\begin{equation}
 \sum_{x,y}|K_{xy}(J^\flat)|+
 \sum_{x,y}|K_{xy}(\widehat J)|
       \ll_\varepsilon N^{-4/5+\varepsilon}.
 \label{p1:limits:erasure-covariance}
\end{equation}
In particular, the absolute second factorial cumulants of
$W$ and $\widehat W$ have this bound.  On the natural coupling,
\begin{align}
 0\le\beta-r_N&\le\tfrac12\E(\Xi)_2=O(a_N^2),
       \label{p1:limits:edge-occurrence}\\
 \TV(\Law W,\Law\widehat W)&\le\tfrac12\E(\Xi)_2=O(a_N^2),
       \label{p1:limits:one-all-tv}\\
 W_1(\Law W,\Law\widehat W)&\le\E(\Xi)_2=O(a_N^2).
       \label{p1:limits:one-all-wasserstein}
\end{align}
Here $a_N\asymp L N^{-2/3}$ as in the main theorem.
\end{proposition}
\begin{proof}
Extend the one-edge deletion kernel to $\mathbb N^{I_N}$ by
choosing uniformly an edge whose two endpoint counts are positive
and subtracting one at each endpoint; if no such edge exists,
leave the configuration unchanged.  It agrees with the binary
deletion kernel.  Apply it to the real and compound fields.  The unweighted comparison error cannot increase.
In the target representation, $T=\sum_eV_e\sim\Pois(\beta)$.
If $T=0$, deletion leaves $U$ unchanged unless two occupied
single sites form an edge.  This event has probability at most
$|E_N|p^2=O(B/N)$.  If $T=1$, let $e$ be its edge.
Deleting $e$ recovers $U$ exactly.  A different edge can be
active only if its endpoints are occupied singles, or if it
shares one endpoint with $e$ and its other endpoint is a
single.  The corresponding bounds are $|E_N|p^2$ and
$O(Bp)$, uniformly in $e$.  Finally $\P(T\ge2)=O(\beta^2)$.
This couples the once-deleted target to $U$ with error
$O(B/N+\beta^2)$.  For isolation, set all coordinates incident
to an active edge to zero; the same argument also excludes
singles at the two endpoints of $e$, at cost $2p$ when
$T=1$.  This proves \eqref{p1:limits:erasure-tv}.

Every deleted count is bounded by its original total count.
The weighted argument at the end of
Lemma~\ref{p1:limits:vector-reference} therefore applies also
to the two deleted laws and $\mathsf U_N$.  With weight
$(1+|\boldsymbol z|)^2$ it controls, in the sum of absolute
entries, their first and second factorial moment measures.
The total intensities are bounded, so it controls the tensor
squares of their means as well.  The independent Poisson field
has zero factorial covariance matrix.  This proves
\eqref{p1:limits:erasure-covariance}.

For the finer intensity assertion, let $B=J-J^\flat$.
Nonnegativity gives the exact identity
\[
 \sum_x(\iota_x-\E B_x)=2(\beta-r_N),\qquad
 0\le\E B_x\le \iota_x.
\]
For isolation, put $D_x=\sum_{y:\{x,y\}\in E_N}J_y$.
Then
\[
 \E\widehat J_x=\E J_x-\iota_x+
       \E[J_x(D_x-1)_+],\qquad
 \sum_x\E[J_x(D_x-1)_+]\le2(\beta-r_N).
\]
The last inequality holds realization by realization: if there
are active edges, at least two vertices are incident to them,
whereas the sum of their active degrees is $2\Xi$.
Now use $\sum_x|\E J_x-p|=N^{-1+o(1)}$ and the positive
two-edge bound in Lemma~\ref{p1:arith:twogoodpositive}.

Finally $0\le k-\ind_{\{k\ge1\}}\le(k)_2/2$ for integers
$k\ge0$.  Also $W=\widehat W$ when $\Xi\le1$, and
\[
 0\le W-\widehat W\le2(\Xi-\ind_{C_N})\le(\Xi)_2.
\]
These inequalities prove
\eqref{p1:limits:edge-occurrence}--\eqref{p1:limits:one-all-wasserstein}.
\end{proof}

The same covariance conclusion applies to the positive cubic
count used to approximate isolation.  With the active degrees
$D_x$ from the proof, put
\[
 Y_x^{[3]}=J_x\left(1-D_x+\binom{D_x}{2}\right),\qquad
 Y^{[3]}=\sum_xY_x^{[3]},\qquad
 \Sigma=\sum_xJ_x\binom{D_x}{2}.
\]
For an integer $d\ge1$, $1-d+\binom d2=\binom{d-1}{2}\ge0$,
whereas its value at zero is one.  Thus
$0\le\widehat J_x\le Y_x^{[3]}$, and the fields agree when
$\Xi\le1$.  More precisely, if
$O=\sum_xJ_x(D_x-1)_+$, then
\[
 Y^{[3]}-\widehat W=\Sigma-O\ge0,\qquad
 \Sigma\le\tfrac12(\Xi)_2,\qquad
 \E Y^{[3]}=m-2\beta+\E\Sigma.
\]
The inequality for $\Sigma$ counts pairs of distinct good
edges sharing a distinguished endpoint; in a simple graph
each unordered pair can share at most one endpoint.
It follows that both the total-variation and the
$W_1$ distances from $Y^{[3]}$ to $\widehat W$ are
$O(a_N^2)$.  Since the deterministic degree is $O(B)$,
$Y^{[3]}\ll B^2Z$.  The fixed moments of $Z$, H\"older's
inequality on $\{\Xi\ge2\}$, and
Lemma~\ref{p1:arith:twogoodpositive} give, for every fixed
polynomial weight, a difference
$O_\varepsilon(N^{-4/3+\varepsilon})$ between the two laws.
In particular,
\begin{equation}
 |\Var(Y^{[3]})-\E Y^{[3]}|
       \ll_\varepsilon N^{-4/5+\varepsilon}.
 \label{p1:limits:cubic-cumulant}
\end{equation}
This deduction needs no asymptotic for the rarer configurations
where the cubic and isolated counts differ.

Let $P_0=\Law(Z\mid C_N^c)$, $\mu_0=\E[Z\mid C_N^c]$,
and $m=\E Z$.  The preceding proof also gives, for every fixed
polynomial weight, a comparison of $P_0$ with
$\Pois(\lambda-\theta)$ at order $N^{-4/5+\varepsilon}$.
Indeed absence of an observed edge in the compound field forces
$T=0$; the remaining conditional law is $U$ conditioned on
having no accidental edge.  Removing this conditioning costs
$O(B/N)$, and $\P(C_N^c)=1-o(1)$ permits conditioning the
comparison in Lemma~\ref{p1:limits:vector-reference}.
Consequently
\begin{equation}
 2r_N=\theta+O(a_N^2),\qquad
 \E W=m-2r_N=\lambda-\theta+O_\varepsilon(N^{-1+\varepsilon}),\qquad
 \mu_0=\lambda-\theta+O_\varepsilon(N^{-4/5+\varepsilon}).
 \label{p1:limits:background-means}
\end{equation}
The identity for the mean of $W$ is exact before its last
estimate.  One may replace $\theta$ by the two-channel expression
of Theorem~\ref{p1:arith:pairs}, or by the actual second
factorial cumulant, at cost $O_\varepsilon(N^{-4/5+\varepsilon})$.

For clarity, the corresponding localization statements are
\begin{align}
 &\big|\TV(\Law Z,H_{\lambda,\theta})
       -(1-r_N)\TV(P_0,\Pois(\mu_0))\big|
       \ll_\varepsilon N^{-4/5+\varepsilon},
       \label{p1:limits:background-localization}\\
 &\big|\TV(\Law Z,H_{\lambda,\theta})
       -\TV(\Law W,\Pois(\E W))\big|
       \ll_\varepsilon N^{-4/5+\varepsilon},
       \label{p1:limits:erased-localization}\\
 &\big|\TV(\Law((J_x)_x),\mathsf H_N)
       -\TV(\Law J^\flat,\mathsf U_N)\big|
       \ll_\varepsilon N^{-4/5+\varepsilon}.
       \label{p1:limits:spatial-localization}
\end{align}
They follow by the triangle inequality from the comparisons just
proved; the changes of Poisson means cost at most their absolute
differences.  Thus the full stopped expansion supplies, in
particular, the first-order reduction to the background.
Small second cumulants alone would not supply any of these
distributional comparisons.

\subsection{Signed Palm identities and weighted recurrence errors}

It is useful to state precisely which information a scalar
Stein comparison retains.  For a probability $P$ on
$\mathbb N$ with finite first moment, put $P_k=P\{k\}$ and
$P_k=0$ for $k<0$, and define
\begin{equation}
 r_k(P;\mu,\theta)
       =kP_k-(\mu-\theta)P_{k-1}-\theta P_{k-2},\qquad
 \mathcal R(P;\mu,\theta)=\sum_{k\ge0}|r_k(P;\mu,\theta)|.
 \label{p1:limits:raw-residual}
\end{equation}

\begin{lemma}[A scalar recurrence criterion]
\label{p1:limits:stein}
If $0\le\theta<\min(\mu,1/2)$, then
\begin{equation}
 \TV(P,H_{\mu,\theta})
       \le\frac{\mathcal R(P;\mu,\theta)}{1-2\theta}.
 \label{p1:limits:stein-bound}
\end{equation}
If $m=\E X>0$, $\kappa=\Var X-\E X$ and
$0\le\kappa<\min(m,1/2)$, the size-biased variable $X^s$
satisfies the exact identity
\begin{equation}
 \mathcal R(\Law X;m,\kappa)
 =2m\TV\left(\Law X^s,
       (1-\kappa/m)\Law(X+1)+(\kappa/m)\Law(X+2)\right).
 \label{p1:limits:size-bias}
\end{equation}
\end{lemma}
\begin{proof}
For an indicator test $h$, a Poisson Stein solution satisfies
$\mu f(k+1)-kf(k)=h(k)-\Pi_\mu h$ and may be chosen with
$\|f\|_\infty\le1$, $\|\Delta f\|_\infty\le1$.
One proof uses the immigration--death process at time $t$,
whose law from $k$ is
$\mathrm{Bin}(k,e^{-t})+\Pois(\mu(1-e^{-t}))$.
One extra initial individual survives with probability $e^{-t}$,
and two survive with probability $e^{-2t}$.  Integrating these
couplings in the Poisson equation gives respectively the bounds
$\int_0^\infty e^{-t}dt=1$ and
$\int_0^\infty2e^{-2t}dt=1$; set $f(0)=f(1)$.

The operator
$A_\theta f(k)=\mu f(k+1)+\theta\Delta f(k+1)-kf(k)$
has mean zero under $H_{\mu,\theta}$, by applying the elementary
Poisson identity separately to its single and double counts.
Changing indices gives
$P A_\theta f=-\sum_kr_kf(k)$.
The difference of the expectations of $\Delta f(k+1)$ under
$P$ and $H_{\mu,\theta}$ is at most
$2\TV(P,H_{\mu,\theta})$.  Taking the supremum over $h$
and absorbing this perturbation proves
\eqref{p1:limits:stein-bound}.  The two measures in
\eqref{p1:limits:size-bias} differ at $k$ by
$r_k(\Law X;m,\kappa)/m$; total variation is half their
$\ell^1$ distance.
\end{proof}

The constant in \eqref{p1:limits:stein-bound} is not claimed
optimal.  The criterion is sufficient, and the raw norm is not
bounded by total variation alone.  Parameter changes satisfy
the elementary bound
\begin{equation}
 |\mathcal R(P;\mu,\theta)-\mathcal R(P;\mu',\theta')|
       \le|\mu-\mu'|+2|\theta-\theta'|.
 \label{p1:limits:residual-parameters}
\end{equation}
When the first three moments exist and $\mu=m$, $\theta=\kappa$,
the signed residual has
$\sum_kr_k=\sum_kkr_k=0$ and
$\sum_kk^2r_k=\kappa_3^{(F)}(X)$.
The last identity tests it by an unbounded function and is not
an estimate for its norm.

\begin{proposition}[The signed Palm measure]
\label{p1:limits:signed-palm}
For any finite family of indicators, let $Z=\sum_xJ_x$,
$p_x=\E J_x$, $m=\E Z$, $P=\Law Z$, and $T_x=Z-J_x$.
Omit sources with $p_x=0$, and define
\[
 \mathcal D=\sum_xp_x
       \{\Law(T_x\mid J_x=1)-\Law(T_x)\}.
\]
Then
\begin{equation}
 \big|\mathcal R(P;m,\theta)
       -\|\mathcal D-\theta\Delta*P\|_1\big|
       \le2\sum_xp_x^2.
 \label{p1:limits:signed-palm-bound}
\end{equation}
Writing $q_A=\E\prod_{x\in A}J_x$, the coefficient of
$(z-1)^j$ in its generating polynomial is exactly
\begin{equation}
 d_j=\sum_{|A|=j}\sum_{x\notin A}(q_{A\cup\{x\}}-p_xq_A).
 \label{p1:limits:signed-palm-coefficients}
\end{equation}
If all $n$ marginals equal $p$ and
$B_j=\E\binom Zj$, this becomes
$d_j=(j+1)B_{j+1}-p(n-j)B_j$.
\end{proposition}
\begin{proof}
Let $SP$ translate a measure by one and let $(KP)_k=kP_k$.
The size-bias identity is
$KP=S\sum_xp_x\Law(T_x\mid J_x=1)$.
Consequently the residual in \eqref{p1:limits:raw-residual}
is the translate of
\[
 \mathcal D+\mathcal B-\theta\Delta*P,\qquad
 \mathcal B=\sum_xp_x\{\Law(T_x)-P\}.
\]
In the natural coupling, $T_x$ and $Z$ differ only when
$J_x=1$, so $\|\mathcal B\|_1\le2\sum_xp_x^2$.
Translation preserves $\ell^1$ norm, and the reverse triangle
inequality proves \eqref{p1:limits:signed-palm-bound}.
Expanding $\prod_{y\ne x}(1+(z-1)J_y)$ before summing over
$x$ gives \eqref{p1:limits:signed-palm-coefficients}.
\end{proof}

The order of summation matters: both the sum over intersections
of a fixed size and the sum over sources occur before the norm.
In this model $p_x\le2p$ for large $N$, so the diagonal error is
$O(N^{-1})$.  Taking absolute values separately for each
forced coordinate would remove cancellations present in
$\mathcal D$.

For reference, two sufficient uncorrected criteria follow without
any affine assumption.  If
\[
 v_x=\TV(\Law T_x,\Law(T_x\mid J_x=1)),\qquad
 w_x=W_1(\Law T_x,\Law(T_x\mid J_x=1)),
\]
then
\begin{equation}
 \mathcal R(P;m,0)\le2\sum_xp_x(p_x+v_x),\qquad
 \TV(P,\Pois(m))\le\sum_xp_x(p_x+w_x).
 \label{p1:limits:palm-sufficient}
\end{equation}
The first follows from the preceding proposition.  For the
second, couple the two scalar $T_x$ laws optimally and lift
that coupling to the original joint law using conditional
distributions given $T_x$.  The extra difference from $Z$ to
$T_x$ has mean $p_x$.  The Poisson Stein equation with
$\|\Delta f\|_\infty\le1$ gives the displayed bound.

The polynomially weighted comparison already proved also makes
the scalar recurrence comparisons legitimate after deletion.
In particular, with
$\mathcal R_0(P;\mu)=\mathcal R(P;\mu,0)$,
\begin{align}
 &\big|\mathcal R(\Law Z;\lambda,\theta)
       -(1-r_N)\mathcal R_0(P_0;\mu_0)\big|
       \ll_\varepsilon N^{-4/5+\varepsilon},
       \label{p1:limits:background-residual}\\
 &\big|\mathcal R(\Law Z;\lambda,\theta)
       -\mathcal R_0(\Law W;\E W)\big|
       \ll_\varepsilon N^{-4/5+\varepsilon}.
       \label{p1:limits:erased-residual}
\end{align}
Indeed the operator in \eqref{p1:limits:raw-residual} is bounded
from the norm $\sum_k(1+k)|\sigma_k|$ to $\ell^1$, uniformly
for the present parameters; its value on the compound target
is zero.  The analogous statement holds for the Poisson
background.  Lemma~\ref{p1:limits:vector-reference}, the deletion
comparison, and \eqref{p1:limits:residual-parameters} prove
both formulas.  This argument uses fixed moments as well as
the distributional approximation.

\subsection{Removing the remaining active pair dependencies}

The recurrence norm weighted by $1/k$ has a useful stability
property under small changes of a law.  For any probability
$P$ on $\mathbb N$ and $\mu>0$ define
\[
 \mathcal V(P;\mu)=\sum_{k\ge1}\frac{|kP_k-\mu P_{k-1}|}{k}.
\]
It is always finite, even without a first-moment assumption.

\begin{lemma}[A recurrence norm equivalent to total variation]
\label{p1:limits:weighted-residual}
For every $P$ and $\mu>0$,
\begin{equation}
 e^{-\mu}\TV(P,\Pi_\mu)\le\mathcal V(P;\mu)
       \le2(1+\mu)\TV(P,\Pi_\mu).
 \label{p1:limits:weighted-equivalence}
\end{equation}
Moreover, if $0<\mu,\nu\le M$, then
\begin{equation}
 |\mathcal V(P;\mu)-\mathcal V(Q;\nu)|
       \le2(1+M)\TV(P,Q)+|\mu-\nu|.
 \label{p1:limits:weighted-stability}
\end{equation}
\end{lemma}
\begin{proof}
Write $a_k=P_k-\Pi_\mu\{k\}$ and
$b_k=a_k-(\mu/k)a_{k-1}$ for $k\ge1$.
The upper bound follows by summing absolute values.
The forward recurrence gives
\[
 a_k=\frac{\mu^k}{k!}a_0+
           \sum_{j=1}^k\frac{\mu^{k-j}j!}{k!}b_j.
\]
The sum over $k\ge j$ of the coefficient of $b_j$ is at most
$e^\mu$.  Since $\sum_ka_k=0$, it follows first that
$|a_0|\le\sum_j|b_j|$ and then that
$\sum_k|a_k|\le2e^\mu\sum_j|b_j|$.
All series are absolutely convergent by the same coefficient
bound.  This proves the lower estimate.  Applying the upper
operator bound to $P-Q$ and bounding
$\sum_{k\ge1}Q_{k-1}/k\le1$ proves stability.
\end{proof}

Let $E_N^{\rm all}$ be the graph on $G$ containing contacts
at distance exactly $L$ and all separated pairs with nonzero
joint relative kernel.  It contains $E_N$.  Define
\[
 \Xi^{\rm bad}=\sum_{e\in E_N^{\rm all}\setminus E_N}J_e,\qquad
 \Lambda^{\rm loc}=\sum_{x\notin G}J_x,\qquad
 \widetilde J_x=J_x\ind_{\{x\in G\}}
       \ind_{\{\sum_{y:\{x,y\}\in E_N^{\rm all}}J_y=0\}},\qquad
 \widetilde W=\sum_x\widetilde J_x.
\]
There is no degree bound imposed on $E_N^{\rm all}$.

\begin{proposition}[All active pair dependencies can be removed]
\label{p1:limits:all-pairs}
For every $\varepsilon>0$,
\begin{align}
 \tau_N:=\E\Xi^{\rm bad}+\E\Lambda^{\rm loc}
       &\ll_\varepsilon N^{-4/5+\varepsilon},
       \label{p1:limits:bad-mass}\\
 \TV(\Law\widetilde W,\Law\widehat W)&\le\tau_N,\qquad
 0\le\E\widehat W-\E\widetilde W\le2\tau_N.
       \label{p1:limits:all-pairs-cost}
\end{align}
In every realization, distinct retained original spaces are
pairwise direct and their starts are strictly separated.
The values of $\mathcal V$ at the respective actual means differ
by $O_\varepsilon(N^{-4/5+\varepsilon})$.  Thus, for a fixed
$0<\varepsilon<2/15$, the condition
\[
 \mathcal V(\Law\widetilde W;\E\widetilde W)=o(a_N)
\]
is equivalent to the Poisson approximation of the erased
background, and hence to
$\TV(\Law Z,H_{\lambda,\theta})=o(a_N)$.
\end{proposition}
\begin{proof}
A compatible separated dependent pair in $G$ has
$q_e=p^22^{\rho(e)}\ge2p^2$, whence
$q_e\le2|q_e-p^2|$.  The absolute bad-pair estimate therefore
bounds their positive event mass.  Contacts have mass
$N^{-1+o(1)}$ and locally deficient windows have total marginal
mass $N^{-1+o(1)}$.  This proves \eqref{p1:limits:bad-mass}.
Pointwise,
\[
 0\le\widehat W-\widetilde W
       \le2\Xi^{\rm bad}+\Lambda^{\rm loc},\qquad
 \{\widehat W\ne\widetilde W\}
       \subseteq\{\Xi^{\rm bad}+\Lambda^{\rm loc}>0\}.
\]
Markov's inequality gives \eqref{p1:limits:all-pairs-cost}.
Two retained starts cannot be closer than $L$ by incompatibility,
cannot be at distance $L$ by the filter, and cannot have a
nonzero pair kernel by its definition.  The means stay in a
fixed compact subset of $(0,\infty)$, so
Lemma~\ref{p1:limits:weighted-residual} transfers the criterion.
Use \eqref{p1:limits:one-all-tv} and
\eqref{p1:limits:erased-localization} for the final equivalence.
\end{proof}

Pairwise directness here is a property of the retained original
constraint spaces in a given realization.  The filtered
indicators themselves depend on the absence of all neighboring
events.  It does not make these indicators independent, or
exclude primitive dependencies among three or more original
spaces.  The stopped expansion supplies the further information
needed for their distribution.

\subsection{The cost of forcing private prime signs}

A small distance between distributions need not be visible in
the positive cost of a particular conditional coupling.  Here
this distinction can be quantified within the arithmetic model.
We use $\kappa(n)=P^+(\sf(n))$, $P^+(1)=1$.

To state its precise logarithmic scale, let $H_X=\log X$ and
use the following definitions from \cite[Proposition~4.2 and
Appendix~F.3]{pouly2026c}.  For $v\ge e$, let $u(v)\ge1$ solve
$e^{u(v)}/u(v)=v$ and put
\[
 D(v)=vu(v)-\operatorname{Ei}(u(v)),\qquad
 V_X=D(H_X/V_X),\qquad \nu_X=H_X/V_X.
\]
The indicated positive solution is unique for large $X$.
At the precision used below,
\begin{equation}
 2V_X=\sqrt{2H_X(\log H_X+\log\log H_X-\log2-2)}
          +o(\nu_X),\qquad
 \log\log X=o(\nu_X).
 \label{p1:limits:pivot-scale}
\end{equation}

\begin{lemma}[Reciprocal odd-prime pivots]
\label{p1:limits:reciprocal}
As $X\to\infty$,
\begin{equation}
 \sum_{n\le X}\frac1{\kappa(n)}
       =X\exp\{-2V_X+o(\nu_X)\}.
 \label{p1:limits:reciprocal-sum}
\end{equation}
The sum over $N\le n<2N$ is
$N\exp\{-2V_N+o(\nu_N)\}$.
\end{lemma}
\begin{proof}
The asymptotic of Erd\H{o}s--Ivi\'c--Pomerance
\cite[Theorem~1 and equation~(1.3)]{erdosivicpomerance1986}
gives this logarithmic formula for
$S^+(X)=\sum_{n\le X}1/P^+(n)$.
Its first correction is the explicit square root in
\eqref{p1:limits:pivot-scale}; the later terms are
$O(\sqrt{H_X\log H_X}(\log\log H_X)^2/(\log H_X)^2)=o(\nu_X)$.
Since $\kappa(n)\le P^+(n)$, this proves the lower bound.

For the upper bound, write uniquely $n=a^2v$ with $v$
squarefree and use positivity to obtain
\[
 \sum_{n\le X}\frac1{\kappa(n)}
       \le\sum_{a\le\sqrt X}S^+(X/a^2).
\]
The terms $a>e^{3V_X}$ are $O(Xe^{-3V_X})$.
On the remaining range $\log(X/a^2)=H_X-O(V_X)$ and
$\nu_{X/a^2}\sim\nu_X$.  The explicit function
$G(H)=\sqrt{2H(\log H+\log\log H-\log2-2)}$ has
$G'(H)=o(1)$ uniformly there.  The stated classical remainder,
without differentiation, is uniformly $o(\nu_X)$.
Thus
\[
 S^+(X/a^2)\le
 X\exp\{-2V_X+o(\nu_X)\}a^{-2+o(1)}.
\]
These weights have bounded sum.  This also explains the
square-part reduction following Proposition~F.5 in
\cite{pouly2026c}.

For a dyadic lower bound, send each $n\le N$ with
$\kappa(n)>2$ to the unique $2^jn\in[N,2N)$.
The pivot is unchanged, and each image has $O(\log N)$
preimages.  The omitted integers have squarefree part $1$ or
$2$, so number $O(\sqrt N)$.  The factor $\log N$ is absorbed
in $\exp\{o(\nu_N)\}$.  The upper bound follows from the sum
up to $2N$, whose scales agree at this precision.
\end{proof}

Consider a specific forcing for the signed exact-run field.
Let
$E=\lceil3V_N/\log2\rceil$, $B_*=L+E+2$, and
$Y=2B_*$.  Retain starts whose maximal support
$\{x-1,\ldots,x+L+E\}$ has every pivot larger than $Y$.
Let $W_{x,e,\sigma}$ be the indicator of the exact word
\[
 (-\sigma,\underbrace{\sigma,\ldots,\sigma}_{L+e},-\sigma),
 \qquad 0\le e\le E,\quad \sigma\in\{-1,1\}.
\]
Its mean is $\lambda_{x,e,\sigma}=p2^{-e-2}$.
For a source $a=(x,e,\sigma)$, fix every nonpivot prime sign
and set each private pivot of its actual word to make that
word occur.  The resulting full field $W^a$ has exactly
$\Law(W\mid W_a=1)$: conditional on the nonpivots, its raw
word bits were independent and fair and its prescribed
probability was constant.  This is the private-pivot
conditional coupling of \cite[Proposition~F.7]{pouly2026c}.
Its nonnegative coordinate cost is
\[
 \mathcal E_N
       =\sum_a\lambda_a
               \E\|W-(W^a-\boldsymbol e_a)\|_1.
\]

\begin{proposition}[Exact logarithmic cost of this forcing]
\label{p1:limits:private-exact}
For this specified signed forcing,
\begin{equation}
 \mathcal E_N=\exp\{-2V_N+o(\nu_N)\}.
 \label{p1:limits:private-exact-eq}
\end{equation}
Consequently this positive coordinate-by-coordinate cost is
not $O(N^{-\alpha})$ for any fixed $\alpha>0$.
\end{proposition}
\begin{proof}
A target can change only if it contains a multiple of a
source pivot.  Summing multiples, then reciprocals, bounds
the number of directed source--target possibilities by
$O(NB_*^2\sum_{n\le3N}1/\kappa(n)+NB_*^2)$.
The finite Palm ledger
\cite[Lemma~F.2 and equation~(F.6)]{pouly2026c} therefore gives
\[
 \mathcal E_N\ll Np^2+
 p^2\left\{NB_*+
       NB_*^2\sum_{n\le3N}\frac1{\kappa(n)}
       +R^{\rm val}_2(N,B_*)\right\}.
\]
Here $R^{\rm val}_2$ is the sum of $2^{\rho^{\rm val}}-1$
over separated maximal supports.  The full-value estimate
\cite[Proposition~3.24]{pouly2026c} gives
$R^{\rm val}_2(N,B_*)\le N^{5/3+o(1)}$, since
$B_*=\log_2N+o(\log N)$.  The type weights sum to at most
one after the common factor $p$ has been removed.
Lemma~\ref{p1:limits:reciprocal} proves the upper bound.

For a lower bound use only the type $e=0,\sigma=1$, of
mean $p_0=p/4$, and restrict both maximal supports to pivots
larger than $Y_0=(\log N)^5$.  Rankin's inequality with
$s=9/10$ gives
\[
 \#\{n\le3N:\kappa(n)\le Y_0\}
 \le(3N)^{9/10}\zeta(9/5)
                    \prod_{\ell\le Y_0}(1+\ell^{-9/10})
       =N^{9/10+o(1)}.
\]
Thus only $N^{9/10+o(1)}$ starts are excluded.
For each source choose $q=\kappa(x-1)\le e^{3V_N}$
and count targets with $q\mid y-1$.
The dyadic reciprocal estimate shows that these pairs number
$N^2\exp\{-2V_N+o(\nu_N)\}$.

Deleting excluded sources or targets costs at most
$N^{19/10+o(1)}$ pairs; the tail $q>e^{3V_N}$ costs
$O(N^2e^{-3V_N})$, and nearby pairs cost $O(NB_*)$.
The condition $q^2\mid y-1$ removes at most a proportion
$Y_0^{-1}$ of the main count plus $O(N)$.
For any other source pivot $q'$ and target offset, the
Chinese remainder theorem counts at most $N/(qq')+1$
targets.  All these extra intersections cost at most
$O(B_*^2/Y_0)$ times the main count plus $O(NB_*^2)$.
Finally the full-value host estimate
\cite[Proposition~3.7]{pouly2026c} excludes at most
$N^{3/2+o(1)}$ pairs with a nonzero raw relation.
Every removed term is negligible.

For each remaining pair the two raw word vectors $U,V$
are independent and uniform, while the forcing changes
only the first target bit, by the fair bit $U_0+b_0$.
The prescribed target word and its first-bit translate
are disjoint, each of probability $p_0$.  Therefore
$\E|W_{y,0,+}-W^a_{y,0,+}|=p_0$ exactly.
Multiplication by the source weight $p_0$ and summation
over the surviving pairs prove the lower bound in
\eqref{p1:limits:private-exact-eq}.
\end{proof}

The two original words in each lower-bound pair are
independent.  Their positive contribution is caused by
the chosen forcing.  It is therefore not a lower bound
for their dependence, for total variation, or for the
scalar cost after summing all changed coordinates with
their signs.

There is also an obstruction that does not depend on the
choice of an arithmetic anchor.  Let the maximal support
have length $B_*=O_C(\log N)$, and retain starts with at
most one pivot at most $Y=2B_*$.  For $\iota\in\{0,1\}$,
use the unsigned word of length $B_\iota=L+1+\iota$:
the start event when $\iota=0$, and exact excess zero
when $\iota=1$.  Its mean on this retained set is
$q_\iota=p2^{-\iota}$.  Choose an anchor as any deterministic
function of the integer arguments and their factorizations;
if there is a pivot at most $Y$ in the actual word, it must
be the anchor.  Preserve its value and force the other
private pivots to impose the word up to a global sign.
All other prime signs are left fixed.  Let
\[
 \mathcal E_{\iota,N}
 =\sum_xq_\iota\sum_{y\ne x}
       \E|W_{y,\iota}-W^{x,\iota}_{y,\iota}|
\]
be the same-type, off-site part of its coordinate cost.

\begin{proposition}[An obstruction for every such anchor]
\label{p1:limits:all-anchors}
For every fixed $0<a<1/10$, uniformly over these arithmetic
anchor rules,
\begin{equation}
 \mathcal E_{\iota,N}\ge
       (1-o_{C,a}(1))\,2^{-2\iota}\lambda_N^2N^{-a}.
 \label{p1:limits:anchor-lower}
\end{equation}
The same assertion with an additional factor $1/2$ holds
if, instead, the orientation is any fair randomized
function of the source word alone and only its private
pivots are changed.
\end{proposition}
\begin{proof}
Use the sufficient logarithmic-window census proved in the
formalization appendix.  For every fixed $a>0$, with
$h_N=\lfloor\log N/(4\log2)\rfloor$, all but $o_a(N)$ integers
$j\in I_N$ have at least two ordinary $N^a$-smooth integers in
$[j,j+h_N)$.  On the critical band, this interval is contained
in each actual word under consideration; translating the initial
boundary contributes at most one further source.
Thus the actual word contains the required two candidates for all
but $o_a(N)$ sources.  This argument uses only the two-candidate
census, which is proved internally by Gaussian comparison and
localization.  The stronger Dickman relative-density asymptotic
of Matom\"aki--Radziwi{\l}{\l}
\cite[Corollary~6]{matomakiradziwill2016} gives another route to
the same sufficient conclusion.

Restrict again to maximal supports whose every pivot
exceeds $Y_0=(\log N)^5$.  The preceding Rankin bound
loses only $N^{9/10+o(1)}$ starts.  At most one of the
two small pivots can be the anchor, so at each of the
remaining $(1-o_a(1))N$ sources choose a forced pivot
$Y_0<q_x\le N^a$.  There are at least
$(1-o_a(1))N^{2-a}$ pairs with $q_x\mid y-1$.
The clean-pair deletions from the preceding proof apply:
excluded targets cost $N^{19/10+o(1)}=o(N^{2-a})$;
nearby targets, valuations at least two, and other forced
pivots cost $o(N^{2-a})$ as well.  Full raw dependence
has at most $N^{3/2+o(1)}$ hosts.

For a surviving pair, $U,V$ are independent uniform
vectors in $\mathbb F_2^{B_\iota}$.  If the selected
source index is $i$ and the anchor is $j\ne i$, the flip
is $U_i+U_j+b_i+b_j$, a fair bit independent of $V$.
The allowed relative word is the two-point coset
$\{b,b+\boldsymbol1\}$, disjoint from its translate by
the first coordinate vector.  The expected target
indicator change is therefore exactly $q_\iota$.
Weighting by the source mean gives
$q_\iota^2(1-o(1))N^{2-a}$, proving
\eqref{p1:limits:anchor-lower}.

For the final assertion, write the prescribed source as
$b+g_x(U,\xi)\boldsymbol1$, where $\xi$ is independent
auxiliary randomness and $g_x$ is fair for uniform $U$.
This guarantees the exact conditional source law even
after all nonpivot bits are fixed.  For two distinct
small-pivot indices $i,j$, the flips satisfy
\[
 \delta_i+\delta_j=U_i+U_j+b_i+b_j.
\]
The right side is fair, so
$\P(\delta_i=1)+\P(\delta_j=1)\ge1/2$.
Choose, before drawing signs, an index with flip
probability at least $1/4$.  Repeat the clean-pair count,
now excluding every other source pivot from the target.
Since $V$ is independent of $(U,\xi)$, each target
change has expectation at least $2q_\iota/4$.
This gives the additional factor $1/2$.
\end{proof}

For any proposed power $\alpha>0$, choose a fixed
$a<\min(\alpha,1/10)$ in
\eqref{p1:limits:anchor-lower}; then
$N^\alpha\mathcal E_{\iota,N}\to\infty$.
The quantifier fixes $a$ before $N$ and invokes no
uniform short-interval smoothness estimate with a growing
parameter.  The result concerns only forcings that change
the indicated private pivots.  An orientation using the
outside environment, a different full conditional coupling,
or a signed scalar comparison is not covered.

